\documentclass[11pt]{article}
\usepackage[margin=0.9in]{geometry}
\usepackage{amsmath,amssymb,amsthm,booktabs,mathrsfs}
\usepackage[T1]{fontenc}
\usepackage{lmodern,microtype}
\usepackage[colorlinks=true,linkcolor=blue,urlcolor=blue]{hyperref}
\newtheorem{theorem}{Theorem}[section]
\newtheorem{lemma}[theorem]{Lemma}
\newtheorem{proposition}[theorem]{Proposition}
\newtheorem{corollary}[theorem]{Corollary}
\theoremstyle{remark}
\newtheorem{remark}[theorem]{Remark}
\newcommand{\norm}[1]{\left\|#1\right\|}
\title{Toward a four-dimensional Yang--Mills construction\\
\large Notebook 32: determinant composition and finite coupling tuning}
\author{Research working document}
\date{17 October 2026}
\begin{document}
\maketitle
\begin{abstract}
We identify the actual leading quantum correction and the precise
drift input needed for finite ultraviolet tuning. All statements
use the composite ordinary constraint, its pivot coordinates,
the absolute contour penalty and the retained Gaussian of the
preceding construction.

First, the finite classical Maxwell leakage tends uniformly to
zero along a cofinal nested period family, even when the preceding
homogeneous depth varies. The actual finite quantum map is jointly
continuous in its interaction class and its positive real coupling;
the continuity argument requires no continuous choice of small
representatives.

Second, we derive an explicit convergent trace formula for the
Maxwell coefficient of the pivot and Gaussian determinants.
Global color invariance removes the second stationary response
only after the full scalar chain rule is taken. The constraint,
current, gauge and logarithmic-chart curvature terms remain in
the precision derivatives. At a fixed regulator, exact Haar
pushforward composition and finite-dimensional Laplace asymptotics
prove a cocycle for these actual leading determinants. Exact
infinite classical Maxwell preservation then makes their
coefficients additive across scales.

An exact finite gauge and coarea identity also rewrites the
leading correction as vector, based-ghost and coarse-constraint
determinants. The coarse factors remain explicit. A flat
Rayleigh test shows that the based ghost alone has no uniform
gap at the final physical scale, so a reference comparison
must retain the coupled cancellations. The coarse constraint
covariance is identified with a shifted actual effective-Hessian
inverse. For a fixed sufficiently large positive shift, its
Maxwell contribution is uniformly bounded in depth and volume.

Finally, under an explicitly unproved positive averaged drift
hypothesis for the actual finite leading coefficients, a continuous
shooting argument tunes arbitrarily long finite local histories
to one fixed positive weak coupling. The auxiliary clipping is
proved absent on the selected histories. This is a conditional
tuning theorem, not an evaluation of the drift. A uniform positive
block margin in every varying-depth structural limit would suffice:
qualitative compactness transfers that margin to the actual finite
regulators without a convergence rate.

The remaining one-loop task is a bound on the difference between
the actual composite coefficient and an explicitly normalized
reference logarithm, uniform in depth. Mixed and large-field
histories, continuum observables, nontriviality and reconstruction
remain part of the full existence argument.
\end{abstract}

The objective remains a nontrivial four-dimensional Yang--Mills
quantum theory. This notebook continues the actual homogeneous
local construction of Notebooks 15--31. Its finite local density
identities do not by themselves represent the entire Wilson
measure. Results and conditional inputs are distinguished in
the statements below.
